% F085｜真实数学图；代表函数不构成抽象解的唯一性证明。
\begin{center}
\begin{tikzpicture}
\begin{axis}[width=12.4cm,height=5.1cm,axis lines=middle,
 xmin=-1.5,xmax=3.1,ymin=-3.9,ymax=4.1,
 xtick={-1,0,1,2,3},xticklabels=\empty,ytick={-3,-1,0,1,3},xlabel={$t$},ylabel={$f(t)$},
 title={余切曲线 $f(t)=\cot t$},title style={font=\small,yshift=3pt},
 tick label style={font=\small},label style={font=\small},
 ylabel style={at={(axis description cs:0,1)},anchor=south east,xshift=-18pt},
 clip=false,grid=major,grid style={black!8},enlargelimits=false,
 every axis plot/.append style={thick},samples=161]
\addplot[CExplanation,very thick,domain=.27:2.85] {cos(deg(x))/sin(deg(x))};
\addplot[CExplanation,very thick,domain=-1.3:-.27] {cos(deg(x))/sin(deg(x))};
\draw[dashed,black!35] (axis cs:0.5235987755982988,0) |- (axis cs:0,1.7320508075688772);
\draw[dashed,black!35] (axis cs:1.0471975511965976,0) |- (axis cs:0,0.5773502691896258);
\addplot[only marks,mark=*,mark size=2pt,orange!80!black] coordinates {(0.523598776,1.73205081) (1.04719755,0.577350269)};

\node[font=\small,anchor=north,yshift=-3pt] at (axis cs:-1,-3.9) {$-1$};
\node[font=\small,anchor=north,yshift=-3pt] at (axis cs:1,-3.9) {$1$};
\node[font=\small,anchor=north,yshift=-3pt] at (axis cs:2,-3.9) {$2$};
\node[font=\small,anchor=north,yshift=-3pt] at (axis cs:3,-3.9) {$3$};

\end{axis}
\end{tikzpicture}
\par\small 两次输入 $\pi/6$ 得 $\pi/3$，高度从 $\sqrt3$ 变为 $(3-1)/(2\sqrt3)=1/\sqrt3$；$t=0$ 无定义。
\end{center}
